% Section 4: Method
% Paper: "Just a Glimpse: Temporal Wearing as a Design Paradigm for Transient Mixed Reality"

\section{Method}

We adopted a Research through Design (RtD) approach~\cite{zimmerman2007research} to investigate temporal wearing as a design paradigm for mixed reality. RtD is well-suited to this inquiry because our research questions concern not only \textit{whether} temporal wearing produces distinct interaction patterns, but \textit{how} design decisions---body coupling, transition cost, form factor---shape those patterns in situ. Rather than testing a single device against a baseline, we designed a family of prototypes that instantiate different positions in the temporal wearing design space (Section~3.3), deployed a subset in a public setting, and analyzed the resulting interactions through mixed methods. This approach follows the precedent of foundational wearable computing research~\cite{gemperle1998wearability} and recent IMWUT work on novel wearable form factors~\cite{khurana2019detachable}, which establish design principles through iterative prototyping and empirical grounding.

\subsection{Prototypes}

We designed and built six functional MR prototypes spanning three form factors---handheld, neck-hung, and stationary---each optimized for temporal wearing (Table~\ref{tab:prototypes}). All prototypes use optical see-through or video see-through displays to deliver spatially registered augmented content, and all were designed to achieve sub-two-second transition costs (onset + offset). Detailed technical specifications, construction processes, and design rationale for each prototype appear in Appendix~A.

\begin{table*}[t]
\centering
\caption{Overview of six temporal wearing prototypes across three form factors. Transition cost is the measured mean time from initiating donning to first augmented frame (onset) plus mean time from initiating doffing to full return to unmediated perception (offset). Three prototypes (marked with $\dagger$) were selected for the field deployment.}
\label{tab:prototypes}
\begin{tabular}{p{2.8cm} p{1.8cm} p{4.5cm} p{2.2cm} p{2.5cm}}
\toprule
\textbf{Prototype} & \textbf{Form Factor} & \textbf{Key Feature} & \textbf{Transition Cost} & \textbf{Deployment} \\
\midrule
[PROTOTYPE A]$\dagger$ & Handheld & [PLACEHOLDER: e.g., monocular viewer with grip handle] & [PLACEHOLDER] s & Field deployment \\
\addlinespace
[PROTOTYPE B] & Handheld & [PLACEHOLDER: e.g., binocular peek device] & [PLACEHOLDER] s & Lab evaluation \\
\addlinespace
[PROTOTYPE C]$\dagger$ & Neck-hung & [PLACEHOLDER: e.g., lanyard-suspended viewer] & [PLACEHOLDER] s & Field deployment \\
\addlinespace
[PROTOTYPE D] & Neck-hung & [PLACEHOLDER: e.g., pendant-style AR monocle] & [PLACEHOLDER] s & Lab evaluation \\
\addlinespace
[PROTOTYPE E]$\dagger$ & Stationary & [PLACEHOLDER: e.g., plinth-mounted viewfinder] & [PLACEHOLDER] s & Field deployment \\
\addlinespace
[PROTOTYPE F] & Stationary & [PLACEHOLDER: e.g., wall-mounted stereoscope] & [PLACEHOLDER] s & Lab evaluation \\
\bottomrule
\end{tabular}
\end{table*}

The prototypes were designed to systematically vary body coupling while holding content constant: all six present the same augmented content experience, a [PLACEHOLDER: brief description of AR content, e.g., spatially registered overlay revealing hidden layers of an artwork]. This controlled variation allows us to isolate the effects of form factor and transition cost on temporal wearing behavior. Each form factor pair includes one prototype optimized for minimal transition cost and one that explores a different point in the design space (e.g., wider field of view at the cost of slightly longer onset time).

\subsection{Field Deployment Design}

From the six prototypes, we selected three for field deployment---one per form factor---to investigate the following research questions:

\begin{itemize}
    \item[\textbf{RQ1}] How does transition cost affect re-engagement frequency? We hypothesize that prototypes with lower transition costs will invite more frequent glimpse cycles within a single visit, as the reduced overhead makes repeated donning and doffing more natural.
    \item[\textbf{RQ2}] How does body coupling affect social comfort in co-located settings? We hypothesize that looser body coupling (handheld, stationary) will produce higher social comfort scores among both users and bystanders, as these form factors maintain social legibility and minimize the appearance of isolation.
    \item[\textbf{RQ3}] What temporal wearing behaviors emerge in situ? Beyond our design space predictions, we seek to identify the behavioral patterns, social dynamics, and appropriation strategies that arise when people encounter temporal wearing devices in a naturalistic setting.
\end{itemize}

These questions address distinct aspects of temporal wearing: RQ1 targets the mechanical dimension (how physical design shapes interaction frequency), RQ2 targets the social dimension (how form factor mediates social acceptability), and RQ3 targets the emergent dimension (what unanticipated behaviors arise in practice).

\subsection{Deployment Site}

We deployed the three prototypes at [VENUE PLACEHOLDER], a [PLACEHOLDER: type of venue, e.g., contemporary art gallery / science museum / university exhibition space] in [PLACEHOLDER: city, country]. The venue was selected for several properties that make it well-suited to studying temporal wearing:

\begin{itemize}
    \item \textbf{Naturalistic transient encounters}: Visitors arrive voluntarily, move freely through the space, and dwell at individual exhibits for seconds to minutes---matching the temporal wearing paradigm's target context.
    \item \textbf{Social co-presence}: Visitors typically attend in pairs or small groups, enabling observation of shared glimpsing, narration, and hand-off behaviors.
    \item \textbf{Spatial configuration}: The venue provided a [PLACEHOLDER: e.g., single room / corridor / open gallery] of approximately [PLACEHOLDER] m$^2$, allowing all three prototypes to be positioned within sight of one another and within the field of view of our video cameras.
    \item \textbf{Visitor diversity}: The venue attracts a broad demographic, including families, couples, student groups, and solo visitors, providing variation in social context and prior technology experience.
\end{itemize}

Each prototype was placed adjacent to a [PLACEHOLDER: artwork / exhibit / installation], with minimal signage---a small label indicating that visitors could pick up or lean into the device to ``see something hidden.'' This deliberate understatement was designed to test the \textit{attract} phase of the glimpse cycle (Section~3.5): how do visitors discover and approach a temporal wearing device without explicit instruction?

\subsection{Participants}

A total of N=[PLACEHOLDER] visitors interacted with the prototypes during the deployment period of [PLACEHOLDER: duration, e.g., three days / one week]. Participation was naturalistic: visitors encountered the prototypes as part of their normal visit to the venue and were not recruited in advance. We collected demographic information only from the subset of visitors who consented to exit interviews (N=[PLACEHOLDER], described below). Of the interviewed participants, [PLACEHOLDER: age range, gender distribution, prior MR experience summary]. The study was approved by [PLACEHOLDER: ethics board] (protocol \#[PLACEHOLDER]).

\subsection{Deployed Prototypes}

Three prototypes were deployed, one per form factor:

\begin{itemize}
    \item \textbf{Handheld} ([PROTOTYPE A]): [PLACEHOLDER: 1--2 sentence description]. This prototype represents the lowest body coupling and fastest transition cost, analogous to historical lorgnettes and handheld viewers.
    \item \textbf{Neck-hung} ([PROTOTYPE C]): [PLACEHOLDER: 1--2 sentence description]. The lanyard maintains passive body coupling when not in active use, reducing the approach phase and enabling rapid re-engagement.
    \item \textbf{Stationary} ([PROTOTYPE E]): [PLACEHOLDER: 1--2 sentence description]. The fixed position eliminates the need to carry or manage the device, with the user's body approaching the device rather than vice versa.
\end{itemize}

All three prototypes presented the same augmented content to control for content effects. The prototypes were counterbalanced across three positions in the venue, with positions rotated [PLACEHOLDER: daily / every two days] to mitigate location effects.

\subsection{Data Collection}

We employed four complementary data collection methods:

\paragraph{Interaction Logs.}
Each prototype was instrumented to record timestamped don and doff events, operationalized as [PLACEHOLDER: e.g., IMU-based lift detection for handheld, proximity sensor activation for stationary, accelerometer threshold for neck-hung]. Logged events included: timestamp, event type (don/doff), duration of each glimpse episode, and inter-episode interval (time between consecutive doff and don events). These logs provided quantitative measures of transition timing, glimpse duration, re-engagement frequency, and total interaction time per visitor.

\paragraph{Video Observation.}
Two cameras were positioned to capture the deployment area from complementary angles: one overview camera capturing the full spatial context (visitor trajectories, social groupings, approach behaviors) and one close-up camera focused on the prototype interaction zone (donning gestures, facial expressions, device manipulation). Video was recorded continuously during venue opening hours, yielding approximately [PLACEHOLDER] hours of footage. Video data was used for behavioral coding and provided context for interpreting interaction log patterns.

\paragraph{Exit Interviews.}
We conducted semi-structured exit interviews with N=[PLACEHOLDER: 15--20] visitors, selected through purposive sampling to ensure variation in age, group composition, and observed interaction patterns (e.g., single-glimpse visitors, multi-cycle visitors, visitors who engaged in hand-off behaviors). Interviews lasted 10--15 minutes and covered: initial discovery and approach, the experience of donning and doffing, awareness of bystanders during use, decisions to re-engage or move on, and comparisons between prototypes for visitors who tried multiple devices. The full interview protocol appears in Appendix~B.

\paragraph{Social Comfort Survey.}
Immediately following their interaction, a subset of visitors (N=[PLACEHOLDER]) completed a brief survey assessing social comfort during device use and willingness to re-engage. The survey used 7-point Likert scales adapted from established social acceptability instruments~\cite{koelle2020social, rico2010gestures} and included items on: felt self-consciousness, perceived bystander reactions, comfort with the physical act of donning/doffing, and likelihood of using the device again. Bystander comfort was assessed through a parallel survey administered to [PLACEHOLDER: N] companions or nearby visitors who observed but did not use the device. The full survey instrument appears in Appendix~C.

\subsection{Analysis}

We employed a mixed-methods analysis approach combining quantitative interaction metrics with qualitative behavioral and interview analysis.

\paragraph{Interaction Metrics.}
Descriptive statistics (mean, median, standard deviation, range) were computed for onset time, offset time, glimpse duration, number of glimpse cycles per visitor, and inter-episode interval, broken down by prototype. To address RQ1, we computed Spearman rank correlations between transition cost and re-engagement frequency across prototypes. Non-parametric tests (Kruskal-Wallis with Dunn's post-hoc comparisons) were used to compare metrics across form factors, given the expected non-normal distribution of timing data.

\paragraph{Video Coding.}
Video data was coded using a scheme adapted from Brignull and Rogers' framework for public display interaction~\cite{brignull2003enticing}, extended to capture temporal wearing-specific behaviors. Coded categories included: discovery behavior (how visitors first noticed the device), approach trajectory, donning gesture, glimpse behavior (scanning, fixating, searching), doffing gesture, post-doff behavior (social sharing, verbal narration, return), and bystander reactions. The coding scheme was developed iteratively: two researchers independently coded [PLACEHOLDER: 20\%] of the video data, discussed disagreements, refined category definitions, and achieved a Cohen's $\kappa$ of [PLACEHOLDER] before coding the remaining footage independently. The full coding scheme appears in Appendix~D.

\paragraph{Thematic Analysis.}
Interview transcripts were analyzed using reflexive thematic analysis following Braun and Clarke's six-phase approach~\cite{braun2006thematic}: familiarization through repeated reading, systematic coding, theme generation, theme review, theme definition, and report writing. We adopted an inductive-deductive orientation: initial codes were generated from the data without reference to the temporal wearing framework, then organized into candidate themes that were iteratively refined in dialogue with the framework. This approach balances sensitivity to participants' lived experience with the theoretical commitments of Research through Design. Two researchers conducted the analysis collaboratively, meeting regularly to discuss coding decisions and theme development.
